% Focused AMGX versus polynomial-preconditioner comparison.
\providecommand{\ADRResultsPath}{docs/research/solver_studies/adr_scaling_2026_09_17}
\providecommand{\ADRResultsFigurePath}{run_outputs/solver_studies/adr_scaling_2026_09_17/figures}

\subsection{Focused AMGX--polynomial comparison}
\label{sec:adr-focused-comparison}

Three square-domain problems sharpen the comparison between AMGX and
polynomial-preconditioned GMRES. On $\Omega=[-1,1]^2$ with exact Dirichlet data:
\begin{equation}
 \begin{aligned}
 -\Delta u+(1+x,-\tfrac12+y)^T\cdot\nabla u+\tfrac12u&=f,
 &u_\star&=1+x^2+2y^2 \quad\text{(variable velocity)},
 \end{aligned}
\end{equation}
\begin{equation}
 \begin{aligned}
 -10^{-3}\Delta u+(1,\tfrac12)^T\cdot\nabla u+u&=f,
 &u_\star&=\sin(\pi x)\sin(\pi y) \quad\text{(transport)},
 \end{aligned}
\end{equation}
\begin{equation}
 \begin{aligned}
 -\Delta u+(1,\tfrac12)^T\cdot\nabla u+u&=f,
 &u_\star&=\sin(\pi x)\sin(\pi y) \quad\text{(high diffusion)}.
 \end{aligned}
\end{equation}
The endpoint comparison uses $999\,940$ trace unknowns at $p=3$ and
$1\,003\,625$ at $p=4$.  The scaling comparison repeats
$n=32,64,128,223$ at $p=6$ and $p=1,2,3,4,6$ at $n=64$.

Every GPU solve uses FP64 arithmetic, a zero initial guess, internal tolerance
$10^{-11}$, physical-residual acceptance at $10^{-10}$, and a 1,000-iteration
limit.  After one warmup setup, three fresh setups each perform two solves.
Fresh time is the median setup-plus-first-solve time; reused time is the mean
second-solve time.  The scaling curves freeze one qualified configuration per
method and problem.  AMGX uses natural $(p+1)\times(p+1)$ BSR blocks; its AMG
variants use block Jacobi for the two diffusive problems and multicolor DILU
for transport.  The polynomial families use standalone polynomial preconditioning, block Jacobi, and block additive Schwarz under
GMRES.  All three problems include $p$MG--AMG/GMRES; the transport problem
also includes direct BSR-DILU/BiCGSTAB. The $p$MG--AMG curves select the fastest
passing available policy independently for fresh and reused timing.

\begin{figure}[!htbp]
 \centering
 \includegraphics[width=\linewidth]{\ADRResultsFigurePath/comparison/million_dof_comparison.pdf}
 \caption{All configurations attempted near one million trace unknowns.
 Filled and open symbols denote $p=3$ and $p=4$; crosses denote failed
 acceptance, and blank positions were not selected for that problem.  Each
 point is the best recorded fresh or reused time under the common protocol.
 Jac. abbreviates block Jacobi; LU labels give the MKL thread count.}
 \label{fig:adr-focused-million}
\end{figure}

At both endpoint degrees, a passing AMGX AMG configuration gives the fastest
fresh solve in every problem.  At $p=4$, its fresh-time advantage over polynomially preconditioned ASM
is $6.3\times$ for variable velocity (0.548 versus 3.443~s), $2.7\times$ for
transport (0.703 versus 1.924~s), and $5.6\times$ for high diffusion
(0.580 versus 3.252~s).  With setup reused, the corresponding AMGX advantages
are $13.0\times$, $2.9\times$, and $9.8\times$.  The $p$MG--AMG hierarchy gives the
overall fastest reused $p=4$ variable-velocity solve, 0.126~s, but its fresh
time is 1.234~s.  Thus setup amortization changes the best method, while it
does not reverse the AMGX--polynomially preconditioned ASM comparison at these endpoints.

\subsubsection{Variable-velocity crossover}
The spatially varying drift tests coefficient variation without oscillatory
geometry.  Direct DILU fails acceptance for both outer methods and degrees;
polynomially preconditioned BJ also narrowly fails at $p=3$.  Every multilevel AMGX configuration,
polynomially preconditioned ASM, polynomial preconditioning, and $p$MG--AMG passes.  The fastest AMGX fresh times are 0.419 and
0.548~s at $p=3,4$, compared with 3.129 and 3.443~s for polynomially preconditioned ASM.

Mesh refinement exposes the crossover.  At 21k unknowns, polynomially preconditioned BJ gives the
fastest fresh result (37.6~ms), and polynomially preconditioned ASM and $p$MG--AMG tie within 0.3~ms for
the fastest reused solve.  At 1.04 million unknowns, AMG/FGMRES takes
0.669~s fresh and 0.218~s reused, versus 3.580 and 2.568~s for polynomially preconditioned ASM.
The $p$MG--AMG hierarchy remains the fastest reused method at 0.136~s.  Standalone
polynomial preconditioning reaches the iteration limit at the two finest meshes.

\begin{figure}[!htbp]
 \centering
 \includegraphics[width=\linewidth]{\ADRResultsFigurePath/comparison/h_scaling.pdf}
 \caption{Mesh scaling at $p=6$ for the three focused problems.  Top: setup
 plus first solve; bottom: repeated solve with setup reused.  Every solver family
 is shown; $p$MG--AMG is a metric-specific policy envelope, and crosses at the top edge mark failed acceptance.
 The largest mesh has 99,458 triangles and 1,041,187 trace unknowns.}
 \label{fig:adr-focused-h}
\end{figure}

\subsubsection{Transport-dominated crossover}
Small and moderate systems favor local preconditioners.  Across the fixed
$n=64$ degree sweep, polynomial preconditioning or polynomially preconditioned BJ is always fastest fresh among GPU methods, while polynomially preconditioned ASM or
direct DILU is the fastest GPU option reused. The million-unknown endpoint instead favors
AMGX AMG: its fresh times are 0.587 and 0.703~s at $p=3,4$, versus 1.654 and
1.924~s for polynomially preconditioned ASM.  Direct DILU/FGMRES fails at $p=3$ and is much slower at
$p=4$; DILU/BiCGSTAB passes both.

$p$MG--AMG also converges on every transport system, but its additional
iterations make it slower than polynomially preconditioned ASM at the million-unknown endpoints.
The standard policy is faster than robust despite requiring more iterations;
strengthening this hierarchy is not a time-to-solution improvement here.

At the largest $p=6$ mesh, direct DILU/BiCGSTAB has the best fresh time
(0.936~s), narrowly ahead of AMG/FGMRES (1.027~s).  Reusing setup reverses
them: AMG/FGMRES takes 0.574~s, direct DILU 0.634~s, and polynomially preconditioned ASM 0.683~s.
This is the only focused problem in which polynomially preconditioned ASM approaches the large-system
AMGX reused time; its remaining gap is $19\%$.

\begin{figure}[!htbp]
 \centering
 \includegraphics[width=\linewidth]{\ADRResultsFigurePath/comparison/p_scaling.pdf}
 \caption{Polynomial-degree scaling on 8,192 triangles.  Fresh and reused
 times are separated, and all measured solver families are shown, including $p$MG--AMG transport results.
 AMGX uses BSR AMG in every panel; the transport panels additionally show
 direct BSR-DILU.}
 \label{fig:adr-focused-p}
\end{figure}

\subsubsection{High-diffusion trigonometric robustness}
All 16 endpoint configurations pass.  AMGX AMG gives the fastest fresh and
reused times at both degrees: 0.450/0.128~s at $p=3$ and 0.580/0.219~s at
$p=4$, where each pair is fresh/reused.  polynomially preconditioned ASM needs 3.057/1.657~s and
3.252/2.157~s.  $p$MG--AMG is intermediate for a fresh solve and approaches
AMGX when its hierarchy is reused.

At 1.04 million unknowns and $p=6$, AMG/BiCGSTAB remains fastest at
0.704~s fresh and 0.250~s reused.  $p$MG--AMG reaches 0.276~s reused, while
polynomially preconditioned ASM takes 3.494~s; standalone polynomial preconditioning reaches the iteration limit.  The fixed
8,192-triangle degree sweep still favors polynomial preconditioning or polynomially preconditioned BJ for one fresh solve,
showing that the large-system advantage is a scaling result rather than a
uniform ranking across sizes.

\begin{table}[!htbp]
 \centering\small
 \setlength{\tabcolsep}{3.2pt}
 \input{\ADRResultsPath/comparison/largest_table.tex}
 \caption{Selected family configurations on 99,458 triangles at $p=6$
 (1,041,187 trace unknowns).  Reused time is a mean; fresh time is a median.
 $p$MG--AMG selects the best reused-time policy and reports both of its timings.
 Bold entries are the best passing times among the displayed rows.  NC denotes
 failure of the common acceptance contract.}
 \label{tab:adr-focused-largest}
\end{table}

In the original endpoint and scaling studies, 184 of 194 configurations pass; the
largest passing physical relative residual is $1.65\times10^{-11}$.  All four
scaling failures are standalone polynomial preconditioning at the two finest variable-velocity and
high-diffusion meshes.  Block AMG passes every focused scaling point.  The
results therefore support AMGX as the stronger large-system default, with
local polynomial methods preferred for inexpensive fresh solves at small and
moderate size and $p$MG--AMG competitive when a diffusive hierarchy is reused.
